\section{Scope, Method, and Qualifications}

\subsection{Reader, question, and method}

This is a discursive canon-law article for a serious general reader able to follow documentary argument---secondarily for students of canon law, theology, or jurisprudence who need the classical taxonomy of law presented from its actual sources rather than from summaries. Its governing question: what is the classical taxonomy of law---eternal, natural, divine positive, human---where does it come from, and what work does it actually do in the current canon law of the Latin Church? Its claim, argued in epistemic order rather than announced: the taxonomy is not a four-box grid but two asymmetric distinctions (author; mode of institution) whose vocabulary the Latin tradition received from Roman jurisprudence through Isidore, whose theory is Aquinas's treatise on law (natural law as participation of the eternal law; human law derived by conclusion and determination), whose canonical reception begins on the first page of Gratian's \emph{Decretum}, and whose working presence in the 1983 Code is pervasive and precise---fixing the subjects of merely ecclesiastical law (c.~11), the boundaries of custom (cc.~23--28) and dispensation (cc.~85--93), the exemptions from prescription (c.~199), the canonization of civil law (cc.~22, 1290), the structures held to be of divine institution, and the declare/constitute grammar of c.~1075; with the modern magisterial exposition (CCC 1954--1960; \emph{Veritatis splendor} 40--45) teaching the same participation doctrine against autonomist and heteronomist alternatives.

The argument is governed by the repository's articles profile, under both its faith-and-theology and its canon-law rules: exact body of law and edition identified per claim; promulgation and effective dates verified at the acts themselves; translation witnesses named; text distinguished from interpretation, validity from liceity, obligation from permission; and authority levels kept distinct throughout---a seventh-century encyclopedia, a twelfth-century school collection, a thirteenth-century theological treatise, a promulgated code, a catechism, and an encyclical are cited each as what it is. Project synthesis (the asymmetry formulation of section 6, the three-verb reading of section 10, the claims-audit of section 12) is labeled as such in the notes where it occurs.

\begin{threestable}{Source class}{Function in the article}{Governing boundary}
Isidore, \emph{Etymologies} & Book V chapters ii--vi, xxi quoted as the vocabulary the tradition received (section 2). & Quoted from Lindsay's 1911 Oxford edition at retained page images (leaves 192, 193, 196 of the identified scan), read visually; OCR used as discovery aid only; cited as transmitter, not as authority on Roman law.\\
Aquinas, ST I-II qq.~90--97 & The systematic synthesis, worked through question by question with the exact loci quoted (sections 3--5). & Latin from the Corpus Thomisticum web text (presents the Leonine edition; not collated against print); English from the identified Dominican translation as web-delivered; both fetched and hashed 2026-07-25. School theology of the highest common authority, not magisterium.\\
Gratian, \emph{Decretum} DD.~1, 4, 8, 9 & The canonical reception (section 7). & Friedberg 1879 text as transcribed by the MDZ online edition, dated state 2026-07-25; vulgate text, recension questions noted; never promulgated legislation---cited as the medium of the tradition c.~6 §2 invokes.\\
CIC 1983 & The working uses: cc.~6, 11, 17, 19, 22--28, 85--93, 98, 113, 129, 145, 199, 207, 375, 748--751, 1008, 1059, 1075, 1084, 1163, 1249, 1290, 1315, 1399, quoted exactly (sections 8--10). & Latin and English Vatican web deliveries, all pages fetched, hashed, and read 2026-07-25; Book VI in the \emph{Pascite gregem Dei} revision; web deliveries, not AAS pages; delivery defects recorded, not corrected.\\
CCC; \emph{Veritatis splendor}; \emph{Dignitatis humanae} & The modern exposition (section 11; DH 3 at section 11 and c.~748's background). & Dated Vatican web states, English with Latin control for the CCC; VS Latin (AAS) not collated; teaching documents cited at their own authority levels.\\
Contemporary debates & Reported in section 12. & Literature-level report from the disputants' standard self-descriptions; none of the disputed works consulted at source; no adjudication.\\
\end{threestable}

\subsection{Included and excluded scope}

Included: the taxonomy's sources and structure as documented in the named witnesses; its reception in Gratian's opening distinctions and the 1983 Code; custom and dispensation as worked institutes; the CCC and \emph{Veritatis splendor} expositions; the declare/constitute boundary; the classical doctrine of unjust law read exactly; the contemporary disputes at reported level.

Excluded: the Eastern Catholic churches---no CCEO analysis is performed, and nothing here may be applied to Eastern Catholics without one; the 1917 Code except at incidental reported level; Roman law at source (the \emph{Digest} and \emph{Institutes} stand behind Isidore but were not examined); the treatise on the Old and New Law (ST I-II qq.~98--108) beyond q.~91 a.~5; the philosophical literature on natural law outside the reported debates; civil-jurisdiction questions; particular and proper law; liturgical law; and every concrete case---no marriage, penalty, office, obligation, or conscience question of any person is analyzed or resolved here.

\subsection{Material qualifications}

\begin{enumerate}[leftmargin=1.45em,itemsep=0.25em]
\item All Vatican web texts are delivery states, not the \emph{Acta Apostolicae Sedis}; two delivery defects were found and recorded (c.~11 English ``efficient''; c.~1315 Latin \latin{divinamcongrua}), and the possibility of others is not excluded.
\item The Summa is quoted from an electronic text identifying itself as Leonine, uncollated with print; the English translation is a working aid whose renderings occasionally smooth the Latin (noted where material).
\item Gratian is quoted from Friedberg's vulgate text in the MDZ transcription; recension scholarship since Winroth is acknowledged and not engaged; wording differences between Gratian's Isidore and Lindsay's Isidore are transmission history, reported without resolution.
\item The canons inventory (section 8) is bounded by its stated search method; it is not a concordance, and register assignments in the ``register invoked'' column follow the canons' own wording, with disputed classifications flagged as disputed.
\item The claim that the Code nowhere labels each impediment's register, and the reading of cc.~1085, 1141--1150 questions as disputed territory, are carried at the level of the standard commentaries, which were consulted for orientation only.
\item The dating of Gratian (c.~1140) and all biographical framing are common scholarly background carried at reported level.
\item Doctrinal-authority questions (which declarations of divine law are definitive; the theological note of particular ``divine institution'' claims) are identified as further questions the taxonomy does not settle; nothing here assigns a theological note beyond what the quoted texts assign themselves.
\item The article's synthesis claims---the empty fourth square, the three-verb grammar, the reading of dispensation and custom as authorship-bounded---are interpretations of the quoted sources, presented as the article's own and traceable in the notes.
\end{enumerate}

\subsection{Legal Scope and Currentness}

Governing law: the 1983 \emph{Codex Iuris Canonici}, promulgated 25 January 1983 by \emph{Sacrae disciplinae leges} and in force from 27 November 1983, as amended---including Book VI as integrally revised by \emph{Pascite gregem Dei} (23 May 2021, in force 8 December 2021), whose promulgation facts were verified at the constitutions' official texts. Authoritative language: Latin; the canons were verified at the Holy See's Latin web deliveries (one page per book, fetched and hashed 2026-07-25), with the Holy See's English web deliveries (translator unnamed) as identified working translations; where they diverge, the Latin governs and divergences are reported. Jurisdiction and persons: the Latin Church only (c.~1); no CCEO norm is analyzed, and Latin conclusions must not be transferred to Eastern Catholics. Universal law only: no particular, proper, liturgical, or penal-particular law is examined. Material facts assumed: none---the article analyzes norms in the abstract and determines no person's obligations, status, or liability. Amendments: the canons quoted were checked at the current Vatican deliveries as of the as-of date; the article asserts currentness for the quoted canons as delivered, not a completed search of every amending act affecting other canons; known post-1983 amendments material to quoted canons (the 1998 \emph{Ad tuendam fidem} changes to c.~750; the 2021 Book VI revision) are reflected in the texts used. \textbf{As-of date: 25 July 2026.} This is a study aid, not legal or canonical advice; rights, penalties, marriage status, sacramental access, office, and every concrete application belong to the competent ecclesiastical authority or a qualified canonist with the complete facts.

\subsection{Notes, translations, rights, and review}

The numbered Notes carry exact citations, verification dates, evidentiary ceilings, and claim-local qualifications; no indispensable premise lives only in a note. Latin is quoted where the original governs, with English renderings either quoted from the identified translations or supplied as labeled working glosses. Official Holy See texts, the Corpus Thomisticum and New Advent presentations, and the MDZ Gratian transcription are quoted briefly with attribution from registered dated states and remain outside the project's CC~BY~4.0 grant; the retained Isidore page images are reproductions of a public-domain 1911 printing; project-created prose and organization are project content. All online witnesses were fetched, hashed, and read publication-locally on 2026-07-25; source-library records (dated web states, artifacts with hashes, registered passages, and three retained public-domain page images) were created or bound for this article, with publication bindings in this leaf's research records.

This revision received internal argumentative, source-consistency, quotation, rights, and production review by the authoring agent. Independent review---philosophical, historical, theological, and canonical---is outstanding; no imprimatur, nihil obstat, or ecclesiastical approval is claimed, and internal checking is not independent review. The publication language for this work is: source-audited working article.
